\documentclass{article}
\usepackage{amsmath,amssymb}
\usepackage{hyperref}
\usepackage{siunitx}
\title{Mathematical and Physical Description of BLAST N5K Notebook}
\author{Analysis of BLAST_x_N5K.ipynb}
\begin{document}
\maketitle

\section*{Summary}
This document explains the mathematical and physical computations performed in the notebook BLAST_x_N5K.ipynb: (i) precomputation of radial integrals of spherical Bessel functions weighted by Chebyshev basis functions; (ii) assembly of projected kernels and double integrals over comoving distance $\chi$ and ratio $R=\chi_2/\chi_1$; (iii) adding Limber-modeled non-linear corrections; (iv) Chebyshev interpolation onto the N5K $\ell$-grid and computation of the $\Delta\chi^2$ accuracy metric.

\section*{Core mathematical definitions}
- Chebyshev decomposition of the 3D power spectrum:
$$
P(k,\chi_1,\chi_2)\approx\sum_{n=0}^{n_{\max}} c_n(\chi_1,\chi_2)\,T_n(k),
$$
where $T_n(k)$ are Chebyshev basis functions and $c_n(\chi_1,\chi_2)$ are coefficients computed by FFTs in the notebook.

- Precomputed spherical-Bessel integrals (definition used in the code):
$$
\tilde T_{n;\ell}^{AB}(\chi_1,\chi_2)=
\int_{k_{\min}}^{k_{\max}} \mathrm{d}k\; T_n(k)\; j_\ell(k\chi_1)\, j_\ell(k\chi_2)\; k^\beta,
$$
where $\beta$ depends on the tracer pair AB:
- clustering–clustering (gg / `CC`): $\beta=+2$ (factor $k^2$),
- shear–shear (ss / `LL`): $\beta=-2$ (factor $k^{-2}$),
- clustering–shear (gs / `CL`): $\beta=0$.

(Implementation: blast_code/src/projected_matter.jl and the generalized variant blast_code/src/new_funcs.jl.)

- Projected kernel (inner $k$ contraction) producing the quantity $w_\ell^{AB}$:
$$
w_{\ell}^{AB}(\chi_1,\chi_2)=\sum_{n} c_n(\chi_1,\chi_2)\,\tilde T_{n;\ell}^{AB}(\chi_1,\chi_2).
$$
(Implementation: `w_ell_tullio` in blast_code/src/projected_matter.jl.)

- Outer integrals to form angular power spectra. The code evaluates integrals in $\chi$ (Simpson) and $R$ (Clenshaw–Curtis). With a discrete quadrature,
$$
C_\ell^{AB}(i,j)=\mathrm{pref}_\ell \,\Delta\chi\sum_{n,m}\chi_n\,K_{i,j}( \chi_n,R_m)\,w_{\ell}(\chi_n,R_m)\,w_\chi[n]\,w_R[m],
$$
where $K_{i,j}(\chi,R)$ is the combined probe kernel for tomographic bins $(i,j)$ and $\mathrm{pref}_\ell$ encodes $\ell$-dependent factors (spin, normalization). (Implementation: blast_code/src/integrals.jl.)

\section*{Physical interpretation}
- The inner $k$ integral projects the 3D matter power spectrum (via Chebyshev expansion) onto spherical-Bessel kernels; physically this captures radial oscillations and scale dependence for each multipole $\ell$ and for pairs of comoving distances $(\chi_1,\chi_2)$. 
- The outer integrals over $\chi$ and $R$ combine those projected matter responses with survey-specific window functions (number-density and lensing efficiencies) to produce tomographic angular cross- and auto-spectra $C_\ell^{gg}, C_\ell^{gs}, C_\ell^{ss}$. 
- Non-linear small-scale corrections are added using Limber approximations for the non-linear part of $P(k)$ — the notebook computes the non-Limber linear component with the SFB-like method and replaces the small-scale part with Limber-evaluated non-linear $C_\ell$ (see combination in the notebook cells where final_C\ell = computed + C\ell_limber_nl - C\ell_limber).

\section*{Numerical method mapping to code}
- Precompute Chebyshev grid, nodes, and weights:
  - Clenshaw–Curtis nodes and weights: `get_clencurt_grid`, `get_clencurt_weights` in blast_code/src/projected_matter.jl.
  - Chebyshev coefficient generation via FFT plans and `fast_chebcoefs` (notebook uses a plan and `Blast.fast_chebcoefs`).
- Bessel integrals:
  - `bessel_cheb_eval` constructs $T_n(k)$ samples and evaluates spherical Bessel $j_\ell(k\chi)$ on the $k$ grid (blast_code/src/projected_matter.jl).
  - `compute_T̃` executes the inner sum/integral using Clenshaw–Curtis weights and a threaded loop to produce $\tilde T_{n;\ell}$ arrays (blast_code/src/projected_matter.jl).
  - `compute_T̃_general` extends this to include extra external Bessel factors if needed (blast_code/src/new_funcs.jl).
- Contract Chebyshev coefficients with precomputed $\tilde T$:
  - `w_ell_tullio(c, T)` implements the sum $w_\ell^{AB}=\sum_n c_n\,\tilde T_n$ using `@tullio` for performance (blast_code/src/projected_matter.jl).
- Kernel construction from survey inputs:
  - `blast_tutorials.compute_kernels` reads precomputed kernel arrays and forms $K_{i,j}(\chi,R)$ via symmetric combinations appropriate for auto and cross correlations (blast_code/src/blast_tutorials.jl).
- Outer integrals and prefactors:
  - `simpson_weight_array` and `get_clencurt_weights_R_integration` compute quadrature weights for $\chi$ and $R$ integrals (blast_code/src/integrals.jl).
  - `get_ell_prefactor` returns the $\ell$-dependent normalization for probe combinations (galaxy/ shear / CMB lensing) (blast_code/src/integrals.jl).
  - `compute_Cℓ` performs the double sum using `@tullio` to produce the final $C_\ell$ tensor (blast_code/src/integrals.jl).
- Limber non-linear corrections and blending:
  - The notebook loads Limber-evaluated linear and non-linear $C_\ell$ arrays and forms
  $$
  C_\ell^{\text{final}} = C_\ell^{\text{non-limber,linear}} + C_\ell^{\text{limber,nonlin}} - C_\ell^{\text{limber,linear}},
  $$
  thereby keeping non-Limber linear terms and adding Limber non-linear corrections (notebook cells around Limber combination).
- Interpolation to N5K $\ell$ grid:
  - The notebook multiplies spectra by $\ell^2$, applies Chebyshev interpolation (`chebinterp`), then rescales by dividing by $\ell^2$ to match the N5K $\ell$ points (cell performing `chebinterp`).

\section*{Validation metric}
- The notebook computes the accuracy metric used in N5K:
$$
\Delta\chi^2=\sum_b N_b \operatorname{Tr}\big(\Sigma_b^{-1}\Delta C_b\big)^2,
$$
where $\Sigma_b$ is a bandpower covariance (including shape/shot noise) computed in `blast_tutorials.Σ`, and $N_b$ is the effective number of modes (computed by `get_nmodes_fullsky`) (see blast_code/src/blast_tutorials.jl).

\section*{Practical notes / points to watch}
- Grid consistency: the notebook must recompute $\tilde T$ after any change to the $\chi$, $R$, or $k$ grids. Mismatched shapes (e.g., different R definitions) cause array-shape errors. The notebook contains multiple places where `R` and `nχ` are defined — ensure you run the precomputation cell after changing them.
- Type consistency: many functions are type-parametric on element type `T`. Convert arrays to a consistent floating type (e.g., Float64) before calling `compute_Cℓ`.
- Performance: inner integrals use large $N$ and threaded loops; the code uses `@tturbo`, `@tullio`, and FFT plans to optimize. Adjust `N`, `n_cheb` and threading to trade accuracy vs runtime.
- Limber blending: verify that the transition $\ell$ where Limber replacement happens matches scientific expectations (notebook uses $\ell>200$).

\section*{Key code references}
- Precompute $\tilde T$ and bessel helpers: blast_code/src/projected_matter.jl
- Generalized $\tilde T$ with extra Bessel factors: blast_code/src/new_funcs.jl
- Contract Chebyshev coefficients: blast_code/src/projected_matter.jl
- Kernel construction and covariance helper: blast_code/src/blast_tutorials.jl
- Quadrature, prefactors, and Cℓ assembly: blast_code/src/integrals.jl
- Notebook cells demonstrating pipeline and Limber blending: BLAST_x_N5K.ipynb (cells around precomputation, kernel computation and final combination).

\section*{Suggested filename and compilation}
Save this document as blast_notebook_analysis.tex in your project root and compile with pdflatex/xelatex.

\end{document}
